
\subsection*{Finite dual of $k[x]$}

The following observation might be related to the $q$-holonomicity of the colored Jones polynomials.

Recall
\[
k[x]^\circ := \left\{ f \in k[x]^* \ \middle|\ \ker(f) \supset I \text{ for some cofinite ideal } I \right\}
\]

Since $k[x]$ is PID when $k$ is a field, $I$ is of the form $I = (p(x))$ for some $p(x) \in k[x]$, and hence
\[
\dim_k k[x] / (p(x)) = \deg(p(x)) < \infty.
\]
Hence $f \in k[x]^\circ$ iff $\ker(f)$ contains a nonzero ideal of $k[x]$, equivalently, $f((p(x))) = 0$ for some nonzero polynomial $p(x)$.
If we let
\[
p(x) = c_0 + c_1 x + \cdots + c_r x^r,
\]
and let $a_n := f(x^n), \quad n \geq 0$, then
\[
f((p(x))) = 0 \Longleftrightarrow c_0 a_n + c_1 a_{n+1} + \cdots + c_r a_{n+r} = 0.
\]
In other words,
\[
f(x) \in k[x]^\circ \Longleftrightarrow \text{ a sequence } (a_n) \text{ is linearly recursive over } k.
\]


\subsection*{Factorisability}

Let $H$ be a quasitriangular Hopf algebra with the universal $R$-matrix $R$.
Then the \emph{monodromy element} $Q = R_{21} R$ measures \emph{factorizability} of $H$.

For $X \in H \otimes H$, define a linear map 
\[
T_X : H^* \rightarrow H, \quad f \mapsto (f \otimes 1)(X).
\]
Observe that, for another $Y \in H \otimes H$,
\begin{align*}
    T_X * T_Y (f) & = T_X(f_{(1)})T_Y(f_{(2)}) \\
    & = (f_{(1)} \otimes 1)(X) \cdot (f_{(2)} \otimes 1)(Y) \\
    & = f(X_1 Y_1) X_2 Y_2 \\
    & = T_{XY}.
\end{align*}

In particular, we have
\[
T_Q(f) = T_{R_{21}} * T_{R} = \mu_H  \circ (l^+ \otimes S(l^-)) \circ \Delta_{H^*}
\]
where $l^+, l^-: H^* \rightarrow H$ is defined by $l^+(f) = (1 \otimes f)(R)$, and $l^-(f) = (f \otimes 1)(R^{-1})$.
Note that $S(l^-)(f) = (f \otimes S)(R^{-1}) = (f \otimes 1)(R)$.

If $T_Q$ is surjective, which means that $H$ is factorizable by definition, then we have
\[
H = \mathrm{Im}(T_Q) \subset H_+ S(H_-) \subset H,
\]
where $H_\pm = \mathrm{Im} (l^\pm)$, which implies $H = H_+ S(H_-)$.


\subsection*{Vague intuition on $O_q(SL_2)$}

$O_q(SL_2)$ comes from FRT construction of $2d$ simple $U_q = U_q(\mathfrak{sl}_2)$ representation $V$.
It means that $O_q$ is generated by matrix coefficeitns of $V$ and the braiding on it.
Hence, $O_q(SL_2)$ forms a Hopf algebra in $U_q^\circ$. $(U_q^\circ \cong O_q \# \mathbb{Z}_2?)$
This is compatible with the well-known fact that there is a pairing between $O_q$ and $U_q$.
So maybe we interpret $O_q$ as a Hopf algebra `focused' on $V$.


\subsection*{Note}

Check Kauffman's paper which uses bra-ket notation for tangles.


\subsection*{Relation between then quantum Grassmannian and 4-plats}

Let first take a look for the classical case; $q=1$.

Let $k[x,y]$ be the coordinate ring of the affine plane.
the natural action of $\mathrm{SL}_2$ on the affine plane is interpreted as the coaction on $k[x,y]$.
As the well-known fact in classical invariant theory, we have
\[
    \left( k[x,y]^{\otimes 2} \right)^{\mathrm{SL_2}} \cong k[u]
\]
where $u = x \otimes y - y \otimes x$ which represents a determinant of $2 \times 2$ matrix.
More generally,
\[
    \left( k[x,y]^{\otimes n} \right)^{\mathrm{SL_2}} \cong k[p_{ij}]
\]
where $p_{ij} = x_i \otimes y_j - y_i \otimes x_j$ and $p_{ij}$ satisfies the \emph{Plucker relations}:
\[
    p_{ij} p_{kl} - p_{ik} p_{jl} + p_{il} p_{jk} = 0.
\]

Now let consider the quantized version of the above argument.
The Hopf algebra $k_q[x,y]$ is a braided Hopf algebra in the category of right comodules over $A_q = O_q(\mathrm{SL}_2)$.
As $U_q$-module via the usual pairing between $A_q$ and $U_q$, $k_q[x,y]$ is decomposed into $U_q$ simple modules
\[
    k_q[x,y] \cong \bigoplus_{N=0}^\infty V_N
\]
where $V_N$ is the usual simple $U_q$ representation of $N$-dimensional.
Remember that we have interested in $U_q$-invariants, which is equivalent to $A_q$-coinvariants (if you are sure about non-degeneracy of the duality between $A_q$ and $U_q$).

Now, observe that
\[
    k_q[x,y]^{\underline{\otimes} 2} \cong \bigoplus (V_N \otimes V_M), \quad \left( k_q[x,y]^{\underline{\otimes} 2} \right)^{A_q} \cong \bigoplus (V_N \otimes V_N)^{A_q} = \cong k_q [u]
\]
where $u = x \otimes y - q^{-1} y \otimes x$.
Recall that the coaction is given by
\begin{align*}
    & \begin{pmatrix}
        x \otimes x & x \otimes y & y \otimes x & y \otimes y 
    \end{pmatrix}
    \\
    \longmapsto & \begin{pmatrix}
        x \otimes x & x \otimes y & y \otimes x & y \otimes y
    \end{pmatrix} 
    \begin{pmatrix}
        a^2 & ab & ba & b^2 \\
        ac & ad & bc & bd \\
        ca & cb & da & db \\
        c^2 & cd & dc & d^2
    \end{pmatrix}.
\end{align*}
Therefore, the coaction of $u$ can be computed as follows.
\begin{align*}
    u \mapsto & \begin{pmatrix}
        x \otimes x & x \otimes y & y \otimes x & y \otimes y 
    \end{pmatrix}
    \begin{pmatrix}
        ab - q^{-1} ba \\
        ad - q^{-1} bc \\
        cb - q^{-1} da \\
        cd - q^{-1} dc
    \end{pmatrix} = u.
\end{align*}
Here we use the algebra relations of $A_q$: $ab = q^{-1} ba$, $cd = q^{-1}dc$, $bc = cb$, $ad - da = (q^{-1} - q) bc$, and
\[
\mathrm{det}_q := ad - q^{-1}bc = 1.
\]